\section{Rocket 2D}\label{sec:rocket2d}

\subsection{Frames}\label{sec:frames}
\begin{itemize}
    \item $\mathcal{F}_0$ --- ground/launch frame, origin at the launch point $O^*$.
    \item $\mathcal{F}_g$ --- gravity frame.
    \item $\mathcal{F}_r$ --- rocket (body) frame.
    \item $\mathcal{F}_t$ --- thrust frame.
\end{itemize}
$\mathcal{F}_0$ is fixed at the launch point $O^*$, with $\hat{x}_0$ horizontal (downrange) and
$\hat{z}_0$ down. The state $(x,z,\theta)$ is always integrated in this frame, regardless of which
gravity or ground model is switched on --- it is the one fixed thing in the whole document.
$\mathcal{F}_g$ carries the direction gravity actually points, $\hat{z}_g(\varphi)$, a function of
the downrange arc angle $\varphi$ introduced in \autoref{sec:gravity}. It coincides with
$\mathcal{F}_0$ by default, $\hat{x}_g(0)\equiv\hat{x}_0$ and $\hat{z}_g(0)\equiv\hat{z}_0$, and
only rotates away from it once a non-parallel direction model is enabled --- so $\mathcal{F}_g$
and $\mathcal{F}_0$ are kept notationally separate precisely because ``ground'' and ``the
direction gravity points'' are not always the same thing, even though they agree at the launch
point.

$\mathcal{F}_r$ is right-handed, with $\hat{x}_r$ along the rocket's nose axis (positive forward),
$\hat{z}_r$ completing the pitch plane with it, and $\hat{y}_r$ out of the pitch plane. Pitch
$\theta$ is measured from $\hat{x}_0$ (the ground-frame horizontal), positive nose-up, and the
pitch rate is $\omega=\dot\theta$, the rotation rate about $\hat{y}_r$. In ground-frame
components,
\begin{equation}\label{eq:body_axes_in_ground}
    \hat{x}_r = \cos\theta\,\hat{x}_0 - \sin\theta\,\hat{z}_0, \qquad
    \hat{z}_r = \sin\theta\,\hat{x}_0 + \cos\theta\,\hat{z}_0.
\end{equation}

$\mathcal{F}_t$ carries the actual line of thrust action, $\hat{x}_t$, which need not coincide
with the nose axis: a real nozzle can be canted, or otherwise misaligned, by a small angle
$\delta$ from the body axis --- a fixed property of how the motor is mounted, not of the flight
state. Giving thrust its own frame, rather than folding $\delta$ into every thrust term by hand,
is what lets that distinction be written down once:
\begin{equation}\label{eq:thrust_frame}
    \hat{x}_t = \cos\delta\,\hat{x}_r + \sin\delta\,\hat{z}_r, \qquad
    \hat{z}_t = -\sin\delta\,\hat{x}_r + \cos\delta\,\hat{z}_r, \qquad \hat{y}_t\equiv\hat{y}_r.
\end{equation}
For now $\delta$ is carried as a constant input --- a fixed manufacturing-type misalignment, swept
by hand to see how sensitive the airframe is to it. A commanded, time-varying gimbal angle driven
by a control law is future work (\autoref{sec:backlog}).

\subsection{Forces and moments}\label{sec:forces}
Weight is the main force the rocket has to overcome to increase its altitude:
\begin{equation}
    \vec{W} = mg\hat{z}_g,
\end{equation}
where $m$ is the rocket's mass and $g$ the gravitational acceleration along the local vertical
$\hat{z}_g$. Both factors are models rather than constants, and are chosen independently: the
magnitude $g$ and the direction $\hat{z}_g$ answer different questions --- how strong, and which
way --- and are supplied separately by \autoref{sec:gravity}.

Thrust is generated by ejecting propellant at high velocity, plus a contribution from the pressure
difference across the nozzle exit. With the thrust frame of \autoref{eq:thrust_frame} it is simply
\begin{equation}
    \vec{T} = T\hat{x}_t,
\end{equation}
purely axial in its own frame, with $T$ the thrust magnitude (\autoref{sec:thrust}); the cant
$\delta$ inside $\hat{x}_t$ carries the whole of the misalignment, so this equation never needs to
be rewritten as $\delta$ changes.

Motion through the surrounding air produces an aerodynamic force,
\begin{equation}
    \vec{A}_f = -A_x\hat{x}_r - A_z\hat{z}_r,
\end{equation}
with $A_x,A_z$ the drag and lift magnitudes, defined in the body frame (\autoref{sec:aero}). By
defining a reference point $D$ on the rocket, a pitching moment can also be written down:
\begin{equation}
    \vec{A}_m = -A_m\hat{y}_r,
\end{equation}
with $A_m$ the pitching-moment magnitude, about $\hat{y}_r$.

\subsection{Geometry}\label{sec:geometry}
Each of the three forces acts at its own point on the rocket. Every dimension below is defined in
the rocket frame $\mathcal{F}_r$ and measured from the nose tip $A$.

Weight is defined at the centre of gravity, here called $B$:
\begin{equation}
    \vec{AB} = -l_B\hat{x}_r - u_B\hat{z}_r.
\end{equation}
Thrust is defined at $C$, the point where the (possibly canted) thrust line of $\mathcal{F}_t$
crosses the body:
\begin{equation}
    \vec{AC} = -l_C\hat{x}_r - u_C\hat{z}_r.
\end{equation}
The aerodynamic force is defined at the reference point $D$:
\begin{equation}
    \vec{AD} = -l_D\hat{x}_r.
\end{equation}
The added $u$ dimensions refer to the axis misalignment: in a real rocket, the centre of gravity
and the nozzle can't be strictly on the axis, whether from manufacturing tolerance or from an
as-built mass asymmetry. As $D$ is a reference point rather than a measured one, no misalignment
is defined for it ($u_D\equiv0$); any asymmetry in the external surfaces instead shows up in the
aerodynamic coefficients themselves, e.g.\ in polar plots, not in a lateral offset of $D$. $u_B$
and $u_C$ are carried as ordinary constants alongside $l_B,l_C$ throughout what follows, on the
same footing as the thrust cant $\delta$ --- setting them to zero recovers the on-axis geometry
exactly, and there is no separate ``model'' for them beyond that: like $\delta$, they are inputs
representing an as-built offset, in principle obtained from manufacturing metrology of the actual
airframe rather than derived from anything else in this document.

\subsection{Equations of motion}\label{sec:eom}
\begin{gather}
    m\vec{\ddot{x}} = \sum \vec{F} = \vec{W} + \vec{T} + \vec{A}_f \label{eq:eom_trans}\\
    I\ddot{\theta} = \sum \vec{M}(B) \cdot \hat{y}_r = \left[\vec{BC}\times\vec{T} + \vec{BD}\times\vec{A}_f + \vec{A_m}\right]\cdot \hat{y}_r \label{eq:eom_rot}
\end{gather}
Moments are taken about the centre of gravity $B$, so weight contributes none:
$\vec{BB}\times\vec{W}=0$, for whatever direction $\vec{W}$ points. $m$ and $I$ are not constants
once the propellant model of \autoref{sec:mass} is active --- they become $m(t), I(t)$, and the
centre of gravity $B$ itself travels as $l_B(t)$. Reinterpreting $\mathcal{F}_0$ as Earth-fixed
rather than inertial adds one further translational term to \autoref{eq:eom_trans} and nothing at
all to \autoref{eq:eom_rot} (\autoref{sec:noninertial}). \autoref{sec:analysis} expands these two
equations by hand.
